\section{A triangular-array collision criterion}
\label{sec:triangular-array}

Before specializing to open shifts, we isolate the probabilistic input.
This separates the classical Poisson-approximation mechanism from the
dynamical work needed to verify its hypotheses.

For each \(m\), let \(p_m\) be a probability distribution on a finite
set \(E_m\), and define
\[
  Q_s(m):=\sum_{x\in E_m}p_m(x)^s
  \qquad(s\geq1).
\]
Fix an integer \(k\geq2\).  Let
\(X_m^{(1)},\ldots,X_m^{(N_m)}\) be independent with common law \(p_m\),
and put
\[
  W_{m,k}
  :=\sum_{\substack{I\subseteq\{1,\ldots,N_m\}\\|I|=k}}
     \mathbf1_{\{X_m^{(i)}=X_m^{(j)}\text{ for all }i,j\in I\}}.
\]
The expected number of collisions is
\[
  \Lambda_{m,k}:=\mathbb EW_{m,k}
  =\binom{N_m}{k}Q_k(m).
\]

\begin{theorem}[Power-sum criterion with total-variation error]
\label{thm:abstract-collision-criterion}
There is \(C_k<\infty\), depending only on \(k\), such that
\begin{align}
 d_{\rm TV}\bigl(\mathcal L(W_{m,k}),
   \operatorname{Poisson}(\Lambda_{m,k})\bigr)
 &\leq C_k\min\{1,\Lambda_{m,k}^{-1}\}
 \Bigl(
 N_m^{2k-1}Q_k(m)^2 \notag\\
 &\hspace{36mm}+
 \sum_{\ell=1}^{k-1}N_m^{2k-\ell}Q_{2k-\ell}(m)
 \Bigr).
 \label{eq:abstract-tv-bound}
\end{align}
In particular, if
\[
  \Lambda_{m,k}\longrightarrow\tau\in(0,\infty)
\]
and
\begin{equation}
  N_m^{2k-\ell}Q_{2k-\ell}(m)\longrightarrow0
  \qquad(1\leq\ell\leq k-1),
  \label{eq:abstract-overlap-condition}
\end{equation}
then
\[
  W_{m,k}\Rightarrow\operatorname{Poisson}(\tau).
\]
\end{theorem}

\begin{proof}
Index the collision indicators by the \(k\)-subsets of
\(\{1,\ldots,N_m\}\).  Connect two index sets when they intersect.
Indicators attached to disjoint sets depend on disjoint families of
independent samples and are therefore independent, so this graph is a
dependency graph.

Every collision indicator has expectation \(Q_k(m)\).  For a fixed
\(k\)-set \(I\), the number of \(k\)-sets \(J\) satisfying
\(|I\cap J|=\ell\) is
\[
  \binom{k}{\ell}\binom{N_m-k}{k-\ell}.
\]
If \(1\leq\ell\leq k-1\), both indicators equal one exactly when all
\(2k-\ell\) samples indexed by \(I\cup J\) coincide.  Hence
\[
  \mathbb E[J_IJ_J]=Q_{2k-\ell}(m).
\]
The two local-dependence terms in
\citet[Theorem~1]{ArratiaGoldsteinGordon1989} are consequently bounded
by
\begin{align*}
 b_1
 &\leq\binom{N_m}{k}
   \sum_{\ell=1}^k
   \binom{k}{\ell}\binom{N_m-k}{k-\ell}Q_k(m)^2,\\
 b_2
 &=\binom{N_m}{k}
   \sum_{\ell=1}^{k-1}
   \binom{k}{\ell}\binom{N_m-k}{k-\ell}
   Q_{2k-\ell}(m).
\end{align*}
There is no third conditional-dependence term because the dependency
graph is exact.  The Chen--Stein estimate is at most a universal factor
times \(\min\{1,\Lambda_{m,k}^{-1}\}(b_1+b_2)\).  Bounding the
binomial coefficients by powers of \(N_m\) proves
\eqref{eq:abstract-tv-bound}.

If \(\Lambda_{m,k}\to\tau\), then
\[
  N_m^{2k-1}Q_k(m)^2
  =O\bigl(N_m^{-1}(N_m^kQ_k(m))^2\bigr)=o(1).
\]
Together with \eqref{eq:abstract-overlap-condition}, this makes the
right-hand side of \eqref{eq:abstract-tv-bound} tend to zero.
Convergence of the Poisson means completes the proof.
\end{proof}

The next result is the form used for survivor distributions.  It allows
bounded oscillating prefactors, so it does not assume a single precise
power-sum limit.

\begin{corollary}[Exponential-gap criterion]
\label{cor:abstract-exponential-criterion}
Let \(n_m\to\infty\).  Suppose that for every integer
\(s\in\{k,k+1,\ldots,2k-1\}\),
\begin{equation}
  Q_s(m)=a_s(m)e^{-n_mh_s}(1+o(1)),
  \label{eq:abstract-power-asymptotic}
\end{equation}
where
\[
  0<\inf_m a_s(m)\leq\sup_m a_s(m)<\infty.
\]
Assume the strict gaps
\begin{equation}
  h_t>\frac{t}{k}h_k
  \qquad(k<t\leq2k-1).
  \label{eq:abstract-strict-gaps}
\end{equation}
If \(\binom{N_m}{k}Q_k(m)\to\tau\in(0,\infty)\), then
\[
  W_{m,k}\Rightarrow\operatorname{Poisson}(\tau).
\]

If, along a subsequence, \(a_k(m)\to a_k>0\) and
\[
  N_me^{-n_mh_k/k}\longrightarrow\alpha\in(0,\infty),
\]
then
\[
  W_{m,k}\Rightarrow
  \operatorname{Poisson}\!\left(\frac{a_k\alpha^k}{k!}\right).
\]
\end{corollary}

\begin{proof}
The convergence of the mean and the bounds on \(a_k(m)\) imply
\[
  N_m\asymp e^{n_mh_k/k}.
\]
For \(t=2k-\ell\), where \(1\leq\ell\leq k-1\),
\[
  N_m^tQ_t(m)
  =O\!\left(
    \exp\left(n_m\left(\frac{t}{k}h_k-h_t\right)\right)
  \right)=o(1)
\]
by \eqref{eq:abstract-strict-gaps}.  Thus
\eqref{eq:abstract-overlap-condition} holds, and
Theorem~\ref{thm:abstract-collision-criterion} gives the Poisson limit.
On the stated subsequence,
\[
  \binom{N_m}{k}Q_k(m)
  \longrightarrow\frac{a_k\alpha^k}{k!},
\]
which proves the explicit form.
\end{proof}

\begin{remark}[Minimality of the overlap input]
The convergence of the mean alone does not imply a Poisson law.  If
\eqref{eq:abstract-overlap-condition} fails, two nominal
\(k\)-collisions sharing samples may occur on the same scale as isolated
collisions.  The strict inequalities
\eqref{eq:abstract-strict-gaps} are a convenient dynamical way to rule
this out, not a generic consequence of normalization.  In the open
Markov application, Proposition~\ref{prop:strict-renyi-gap} verifies all
of them simultaneously.
\end{remark}

\begin{remark}[Why only finitely many moments are used]
For a fixed multiplicity \(k\), two intersecting \(k\)-sets have union
size between \(k+1\) and \(2k-1\).  The criterion therefore uses only
the power sums \(Q_k,\ldots,Q_{2k-1}\).  It makes no assertion uniform
in a growing collision multiplicity, and no such uniformity is needed
for the results below.
\end{remark}
